\documentclass[12pt,a4paper]{article}
\usepackage[utf8]{inputenc}
\usepackage[T1]{fontenc}
\usepackage[spanish]{babel}
\usepackage{geometry}
\geometry{margin=2.5cm}

\title{Hito 3: Modelado y experimentación}
\author{Plantilla}
\date{Abril 2026}

\begin{document}
\maketitle

\section{Objetivo}
Entrenar, validar y seleccionar el modelo de churn.

\section{Estructura}
\begin{itemize}
  \item Diseño experimental y particionado temporal.
  \item Modelos baseline y candidatos avanzados.
  \item Métricas técnicas y de negocio.
  \item Fairness y robustez temporal.
  \item Trazabilidad de experimentos.
  \item Selección del modelo final.
\end{itemize}

\section{Mermaid pendiente}
\begin{verbatim}
flowchart LR
    A[Features Gold] --> B[Train/Validation/Test]
    B --> C[Experimentos MLflow]
    C --> D[Modelo final]
\end{verbatim}

\section{Checklist}
\begin{itemize}
  \item Protocolo reproducible.
  \item Métricas de negocio alineadas con ROI.
  \item Fairness evaluada.
  \item Modelo final registrado y versionado.
\end{itemize}

\end{document}
